\documentclass[12pt,a4paper]{article}

\usepackage[utf8]{inputenc}
\usepackage[T1]{fontenc}
\usepackage[portuguese]{babel}
\usepackage{geometry}
\geometry{
    a4paper,
    top=3cm,
    bottom=2cm,
    left=3cm,
    right=2cm
}

\usepackage{helvet}
\renewcommand{\familydefault}{\sfdefault}

\usepackage{hyperref}
\usepackage{titlesec}
\usepackage{setspace}

\begin{document}
\onehalfspacing

\begin{center}
    \textbf{Fundação Universidade Federal do ABC}\\
    \textbf{Pró-reitoria de pesquisa}\\
    \small Av. dos Estados, 5001, Santa Terezinha, Santo André/SP, CEP 09210-580\\
    Bloco L, 3º Andar, Fone (11) 3356-7617 | iniciacao@ufabc.edu.br
\end{center}
\vspace{0.8cm}

\noindent \textbf{Relatório (Parcial ou Final) de Iniciação Científica referente ao Edital:} 01/2025 PIC/PIBIC/PIBITI/PIBIC-AF \\[0.5cm]
\noindent \emph{Na capa deverá conter tanto a assinatura do aluno quanto a assinatura do orientador}\\[0.5cm]
\noindent \textbf{Nome do aluno:} \hrulefill \\
\noindent \textbf{Assinatura do aluno:} \hrulefill \\[0.5cm]
\noindent \textbf{Nome do orientador:} \hrulefill \\
\noindent \textbf{Assinatura do orientador:} \hrulefill \\[0.5cm]
\noindent \textbf{Título do projeto:} Problemas de Fluxos em Redes: Teoria, Algoritmos e Implementações \\[0.5cm]
\noindent \textbf{Palavras-chave do projeto:} Otimização Combinatória, Teoria dos Grafos, Fluxos \\[0.5cm]
\noindent \textbf{Área do conhecimento do projeto:} Ciência da Computação, Matemática Discreta \\[0.5cm]
\noindent \textbf{Bolsista:} ( X ) Sim \quad ( \quad ) Não. \\
\textbf{Em caso positivo, informar a modalidade da bolsa:} PIBIC - CNPq \\[1cm]

\begin{center}
    \textbf{Santo André}\\
    \textbf{\today}
\end{center}

\vspace{1cm}

\renewcommand{\contentsname}{Sumário}
\tableofcontents
\newpage

\section{Resumo}
Fluxos em digrafos constituem uma área fundamental da Otimização Combinatória, com amplas aplicações práticas em redes de transporte e alocação de recursos. Este projeto adotou uma abordagem que combina fundamentos teóricos, implementação algorítmica e validação experimental para o estudo de problemas de fluxo máximo e fluxo de custo mínimo. Foram estudados algoritmos bem estabelecidos, além de técnicas de modelagem para problemas correlatos. A validação prática foi realizada através de experimentos computacionais utilizando instâncias de benchmarks reconhecidas.

\section{Introdução}
A Otimização Combinatória é uma área de pesquisa situada na interseção entre Matemática e Ciência da Computação, cujos resultados possuem abrangência em diversas direções. Do ponto de vista matemático, a área mantém forte influência e interseção com a Teoria dos Grafos, uma vez que grafos constituem objetos combinatórios com amplo poder de modelagem e aparecem frequentemente em problemas de otimização. 

Sob a perspectiva computacional, a Otimização Combinatória aborda algoritmos de variados tipos, priorizando a obtenção de algoritmos eficientes que produzam soluções garantidamente ótimas ou com fatores de aproximação conhecidos. Algoritmos não-polinomiais e heurísticas também ocupam lugar de destaque devido à sua relevância prática.

Dentre os tópicos em Otimização Combinatória, problemas de fluxos em digrafos têm um destaque especial. Essa área se concentra no estudo de redes onde recursos, informações ou materiais devem ser transportados de pontos de origem para pontos de destino, respeitando restrições de capacidade e demanda. 

O objetivo geral do projeto é proporcionar uma formação sólida em problemas de fluxos, integrando fundamentos teóricos, implementação algorítmica e validação experimental. Através de uma abordagem metodológica que equilibra teoria e prática, buscou-se desenvolver competências essenciais de pesquisa científica, incluindo a execução de experimentos computacionais e escrita científica.

\section{Fundamentação teórica}
O problema do st-fluxo máximo, onde se busca maximizar a quantidade de fluxo que pode ser enviada de um vértice fonte s para um vértice de destino t, representa um dos problemas mais clássicos da área. Diversos algoritmos como os de Ford-Fulkerson, Edmonds-Karp, Dinic, Goldberg-Tarjan e Orlin foram desenvolvidos para esse problema. 

Já o problema de fluxo de custo mínimo generaliza essa abordagem ao incorporar custos associados ao transporte, buscando minimizar o custo total enquanto satisfaz as demandas da rede. O algoritmo mais clássico para este problema é o \textit{network simplex}, que é essencialmente uma especialização eficiente do algoritmo simplex para este problema. 

Muitos outros problemas combinatórios podem ser modelados como problemas de fluxo, demonstrando a versatilidade e importância central desta área dentro da otimização combinatória. Dentre eles, destacam-se o problema de emparelhamento máximo em grafos bipartidos e o problema de empacotamento de caminhos disjuntos.

\section{Metodologia}

\subsection{Materiais e Métodos}
A metodologia baseou-se em ciclos de estudo através da bibliografia especializada e implementação. Os algoritmos foram implementados em linguagem C++ para garantir eficiência computacional, utilizando estruturas de dados apropriadas para representação de grafos e fluxos. 

Os experimentos foram conduzidos em instâncias padronizadas das coleções DIMACS, permitindo comparação com resultados da literatura. A validação dos resultados seguiu protocolos comuns na área, com análise de desempenho baseada principalmente em métricas de tempo de execução. Este projeto não necessitou de equipamento especializado além de computadores convencionais.

\subsection{Etapas da pesquisa}
As atividades de pesquisa foram estruturadas da seguinte maneira:
\begin{itemize}
    \item \textbf{Atividade 1:} Estudo dos fundamentos de Otimização Combinatória e Teoria dos Grafos.
    \item \textbf{Atividade 2:} Estudo teórico do problema de fluxo máximo e algoritmos clássicos (Ford-Fulkerson, Edmonds-Karp, Dinic e Goldberg-Tarjan).
    \item \textbf{Atividade 3:} Implementação em C++ dos algoritmos de fluxo máximo.
    \item \textbf{Atividade 4 e 5:} Estudo de técnicas de modelagem e implementação de algoritmos para problemas modelados, como o emparelhamento máximo em grafos bipartidos e empacotamento de caminhos disjuntos.
    \item \textbf{Atividade 7 e 8:} Estudo teórico e implementação do problema de fluxo de custo mínimo e algoritmo \textit{network simplex}.
    \item \textbf{Atividades Finais:} Testes experimentais e análise comparativa de desempenho.
\end{itemize}

\section{Resultados e discussão dos resultados}
Apresente detalhadamente os resultados obtidos com as implementações e experimentos computacionais nas instâncias DIMACS.

\section{Conclusões e perspectivas de trabalhos futuros}
Descreva aqui se os objetivos foram alcançados e detalhe quais serão os próximos passos para dar continuidade a esta linha de pesquisa.

\section{Referências}

BONDY, J. A.; MURTY, U. S. R. \textit{Graph Theory}. Editora, Cidade, País, Ano.

COOK, W. et al. \textit{Combinatorial Optimization}. Editora, Cidade, País, Ano.

DIESTEL, R. \textit{Graph Theory}. Editora, Cidade, País, Ano.

EDMONDS, J.; KARP, R. M. Título do artigo. \textit{Nome da revista}, Volume, p. Página, Ano.

FORD, L. R.; FULKERSON, D. R. Título do artigo. \textit{Nome da revista}, Volume, p. Página, Ano.

SCHRIJVER, A. \textit{Combinatorial Optimization}. Editora, Cidade, País, Ano.

\end{document}
